% One-time imported draft. Editable native source; no Markdown regeneration.
\section{Project contract}\label{009-sorting-visualizer-assignment-visualize-insertion-and-selection-sort}

Sort the active prefix of the eight-word values array into signed
ascending order. The input count is zero through eight, and every active
value is an integer from \ensuremath{-}4 through 3. Inactive words may contain
arbitrary 32-bit values and must receive no execution stores, even
temporarily. Preserve count, algorithm, and every supplied guard
throughout execution. Use the Default little-endian layout. The frame is
256 bytes at 0x10010000, followed by frame\_guard at 0x10010100.

Algorithm zero selects adjacent-exchange insertion sort: extend a sorted
prefix one position at a time, exchange an item with its left neighbor
while it is strictly smaller, and stop moving it when equal or greater.
Algorithm one selects selection sort: scan each remaining suffix for its
first minimum, exchange that minimum with the suffix\textquotesingle s
first item if their indices differ, and advance the sorted prefix. Use
strict less-than when updating a minimum. Self-exchanges are not
published. These policies define reproducible visualization states,
including duplicate keys.

\subsection{Public visualization
  interface}\label{009-sorting-visualizer-assignment-public-visualization-interface}

When main runs, reach \texttt{sort\_ready} after rendering the initial
array (phase=0, \texttt{swap\_count}=0), after every completed algorithm
exchange (phase=1, incremented \texttt{swap\_count}), and once at
completion (phase=2, unchanged final \texttt{swap\_count}). At each
entry the array must be a permutation of its initial active values and
the whole frame must match it. Both word stores of an exchange precede
publication. Do not publish a half-exchange. After the final snapshot,
exit normally through done. Empty, singleton, and already sorted inputs
still publish initial and final snapshots, with no extra exchange.

Column x represents active item x. Encode value v as bar height v+5,
from one through eight cells. A cell at screen row y is white 0x00ffffff
exactly when x\textless count and y\ensuremath{\geq}3\ensuremath{-}values{[}x{]}; every other cell is
black zero. This offset encoding makes negative values visible while
preserving signed order; it is not a signed-magnitude plot. Write all 64
frame words on every render, including black and inactive cells.

Provide a callable render function. Its arguments are \texttt{a0}, a
pointer to a word array, and \texttt{a1}, its valid length zero through
eight. It reads only that active prefix, writes the fixed frame, reaches
\texttt{sort\_ready}, and returns through \texttt{ra}. It preserves
\texttt{s0}..\texttt{s7} and \texttt{sp} and changes no source words,
guards, count/algorithm, phase, or \texttt{swap\_count}. Temporary and
argument registers may be changed. The caller owns publication metadata;
an independent render call leaves it unchanged, regardless of its
current values. A caller can pass another valid word-aligned prefix
address, so the function must not rely on main\textquotesingle s
saved-register setup or the values symbol itself. done is the normal
exit continuation used by independent tests.

The guide uses a leaf function: it calls no other function and therefore
does not need to save a nested return address. Equivalent
implementations may use other internal structures if they preserve the
stated contracts. The public labels do not prescribe private loop
labels, scratch registers, or an exact instruction total. A shift-based
insertion implementation needs a different animation contract; this
project explicitly publishes complete adjacent exchanges.

Produce no console output and exit within 100,000 completed
instructions. No host-time animation delay is required. A fast Run can
skip visible repaints; Debug at \texttt{sort\_ready} provides a stable
complete snapshot. The checks observe real arrays and frames rather than
trusting a student-written log of intended swaps.

\subsection{Submit and
  explain}\label{009-sorting-visualizer-assignment-submit-and-explain}

Submit main.asm, a trace of both algorithms on the default input, and a
statement of each sorted-prefix invariant. Explain why order alone does
not prove value preservation, how duplicates affect stability, and how
the rendering call returns without destroying algorithm state. Compare
sorted, reversed, and shuffled eight-item fixtures using both whole-run
and sorting-range instruction/branch counts, with the measured scope
stated.
